\documentclass[11pt,a4paper]{article}
%TC:envir sprechtext 1
%TC:macro \folie [ignore,ignore,ignore,ignore,ignore]
\usepackage[T1]{fontenc}
\usepackage[utf8]{inputenc}
\usepackage[ngerman]{babel}
\usepackage{lmodern,microtype}
\usepackage[margin=21mm,top=20mm,bottom=21mm,headheight=15pt]{geometry}
\usepackage{graphicx,booktabs,tabularx,array,xcolor,fancyhdr,hyperref}
\definecolor{tured}{HTML}{B6302D}
\definecolor{ink}{HTML}{333333}
\definecolor{quiet}{HTML}{626262}
\hypersetup{colorlinks=true,linkcolor=tured,urlcolor=tured,
 pdftitle={Dissertationsvortrag: Erzählfassung für 30 Minuten, Präsentation v1.12},
 pdfauthor={Florian Rzepka},bookmarksopen=false}
\setlength{\parindent}{0pt}
\setlength{\parskip}{7pt}
\setlength{\emergencystretch}{1em}
\linespread{1.06}
\raggedbottom
\pagestyle{fancy}
\fancyhf{}
\fancyhead[L]{\small Florian Rzepka | Dissertationsvortrag}
\fancyhead[R]{\small Erzählfassung zu v1.12}
\fancyfoot[L]{\footnotesize Stand: 1. Oktober 2026}
\fancyfoot[R]{\footnotesize\thepage}
\renewcommand{\headrulewidth}{0.3pt}
\newenvironment{sprechtext}{}{}
\newcommand{\folie}[5]{\clearpage
  \pdfbookmark[1]{Folie #1: #2}{folie#1}
  {\color{tured}\large\bfseries Folie #1: #2\par}
  {\small\color{quiet}Zielzeit #3 min nach Beginn | Für diese Folie: #4 min\par}
  \vspace{2mm}
  \begin{center}\includegraphics[width=\linewidth]{assets/erzaehlfassung_v1_12/slide\ifnum#1<10 00\else 0\fi#1.png}\end{center}
  {\small\color{quiet}\textit{Regiehinweis, nicht mitsprechen:} #5\par}
  \vspace{3mm}
}
\begin{document}
%TC:ignore
\pdfbookmark[0]{Vorbereitung und Zeitplan}{vorbereitung}
{\small\color{tured} DISPUTATION | SPRECHMANUSKRIPT}
\vspace{4mm}

{\LARGE\bfseries Erzählfassung für 30 Minuten\par}
{\large Zur Präsentation \texttt{Diss\_Presentation\_FR\_v1.12.pptx}\par}
\vspace{3mm}

\textbf{Umfang.} Die Fassung folgt den 40 Vortragsfolien der aktuellen
Präsentation. Jede Folie hat einen Sprechtext, eine Abbildung der tatsächlichen
Folie und einen kurzen Regiehinweis. Die Foliennummern beziehen sich auf die
Position in PowerPoint. Der Anhang ab Folie 41 ist für Rückfragen vorgesehen.

\textbf{Zeitplanung.} 29 Minuten für Begrüßung und Vortrag, eine Minute Reserve.
Die ungefähr 3\,428 Sprechwörter ergeben bei 125--130 Wörtern pro Minute
etwa 26--27{,}5 Minuten reine Sprechzeit. Die verbleibende Zeit innerhalb des
29-Minuten-Plans ist für Bildverweise, kurze Pausen und Folienwechsel gedacht.
Die endgültige Dauer ergibt sich erst in der laut gesprochenen Generalprobe.

\renewcommand{\arraystretch}{1.28}
\begin{tabularx}{\linewidth}{@{}>{\raggedright\arraybackslash}Xrr@{}}
\toprule
Abschnitt & Folien & Zielzeit \\
\midrule
Begrüßung, Motivation und Datengrundlage & 1--7 & 0:00--5:05 \\
SOH-Schätzung mit MLP und Lag Sequences & 8--12 & 5:05--8:50 \\
Schätzer, Störungen und Auswertung & 13--21 & 8:50--14:00 \\
Ergebnisse des Robustheitsbenchmarks & 22--26 & 14:00--18:10 \\
Architektur, Pruning und Quantisierung & 27--34 & 18:10--24:15 \\
Genauigkeit, Speicher und Laufzeit & 35--38 & 24:15--27:40 \\
Fazit und Übergang zur Diskussion & 39--40 & 27:40--29:00 \\
\midrule
Reserve & & 29:00--30:00 \\
\bottomrule
\end{tabularx}

\vspace{4mm}
\textbf{Kommission}\par
\begin{tabularx}{\linewidth}{@{}l>{\raggedright\arraybackslash}X@{}}
Vorsitz & Prof. Dr.-Ing. Kai Strunz \\
Gutachterin & Prof. Dr.-Ing. Julia Kowal, TU Berlin \\
Gutachter & Prof. Dr.-Ing. Frank Opferkuch, TH Nürnberg \\
Gutachter & Prof. Dr.-Ing. Weihan Li, RWTH Aachen \\
\end{tabularx}

\textbf{Begrüßung.} Auf Folie 1 ist ein etwa einminütiger Einstieg enthalten.
Der persönliche Bezug zu Professor Opferkuch bleibt bei einem Satzpaar.
Gesprochen reicht die Anrede \glqq Frau Professor Kowal\grqq{} beziehungsweise
\glqq Herr Professor\grqq{}; die vollständigen akademischen Titel stehen hier
zur Orientierung.

\textbf{Lesart.} Nur die Absätze unter dem Regiehinweis sind Sprechtext.
Die Hinweise auf dieser Seite und die Vorbereitung am Ende werden nicht
mitgesprochen. Die vollständigen Gleichungsherleitungen bleiben im Anhang.
%TC:endignore
\folie{1}{Begrüßung und Thema}{0:00--1:00}{1:00}{Nach der Worterteilung beginnen. Die Namen ruhig aussprechen, kurz zur Kommission schauen.}

\begin{sprechtext}

Guten Tag, Herr Professor Strunz, Frau Professor Kowal, Herr Professor Opferkuch und Herr Professor Li. Herzlich willkommen auch allen Gästen.

Vielen Dank an die Kommission für die Begutachtung und die Durchführung der Prüfung. Frau Professor Kowal, Ihnen danke ich besonders für die Betreuung meiner Promotion. Herrn Professor Opferkuch und Herrn Professor Li danke ich auch für die Anreise aus Nürnberg und Aachen.

Herr Professor Opferkuch, Sie haben bereits meine Bachelorarbeit zu Lithium-Ionen-Akkus an der TH Nürnberg betreut. Ich freue mich, dass Sie auch heute wieder dabei sind.

In meiner Dissertation geht es um neuronale Zustandsbestimmung für eingebettete Batteriemanagementsysteme. Ich zeige heute, wie sich Schätzgenauigkeit, Robustheit und Hardwarebedarf gemeinsam betrachten lassen.

\end{sprechtext}

\folie{2}{Warum Zustandsbestimmung direkt im BMS?}{1:00--1:45}{0:45}{Auf die lokalen Speicher und danach auf die Ressourcenbeschränkung zeigen.}

\begin{sprechtext}

Der Ausgangspunkt sind Batteriespeicher, die zuverlässig vor Ort arbeiten müssen. Das betrifft zum Beispiel dezentrale Energiesysteme, wie hier im MG-Farm-Kontext, und mobile Anwendungen außerhalb einer festen Infrastruktur.

Das Batteriemanagementsystem muss wissen, wie viel Ladung noch verfügbar ist und wie stark die Batterie bereits gealtert ist. Diese Größen lassen sich nicht einfach direkt messen. Sie müssen aus Spannung, Strom, Temperatur und der Betriebshistorie geschätzt werden.

Gleichzeitig stehen nur begrenzter Speicher und begrenzte Rechenleistung zur Verfügung. Auch eine dauerhafte Internetverbindung kann man nicht voraussetzen. Deshalb soll die Zustandsbestimmung direkt auf dem Mikrocontroller funktionieren.

\end{sprechtext}

\folie{3}{Drei Forschungsfragen}{1:45--2:35}{0:50}{Die Fragen in der Reihenfolge Modellgestaltung, Zuverlässigkeit und Umsetzung führen.}

\begin{sprechtext}

Daraus ergeben sich drei Forschungsfragen. Erstens: Wie kann ich zeitabhängige Alterung mit einer möglichst einfachen Modellstruktur abbilden? Dabei geht es besonders darum, welche Informationen ich dem Modell als Eingang zur Verfügung stelle.

Zweitens: Wie zuverlässig sind verschiedene Schätzer, wenn Messwerte gestört sind oder der Anfangszustand nicht stimmt? Ein kleiner Fehler unter idealen Bedingungen beantwortet diese Frage noch nicht.

Drittens: Wie lassen sich neuronale Modelle auf einem Mikrocontroller ausführen und dort hinsichtlich Speicher und Laufzeit verbessern?

Diese drei Fragen bilden den Aufbau meines Vortrags: zunächst die Eingabedarstellung, dann der Robustheitsvergleich und schließlich die Optimierung für die Hardware.

\end{sprechtext}

\folie{4}{Einordnung des eigenen Beitrags}{2:35--3:15}{0:40}{Je Spalte einen Beitrag nennen. Die Literaturangaben nicht vorlesen.}

\begin{sprechtext}

Für alle drei Bereiche gibt es bereits etablierte Ansätze. Der Beitrag meiner Arbeit liegt deshalb in den konkreten Untersuchungen und ihrer Verbindung.

Im ersten Teil analysiere ich systematisch die Eingangsmerkmale und die zeitliche Historie für die SOH-Schätzung mit einem MLP. Im zweiten Teil vergleiche ich unterschiedliche Schätzerklassen unter denselben Messstörungen. Im dritten Teil untersuche ich Pruning und Quantisierung für SOC- und SOH-Modelle auf derselben Mikrocontrollerplattform.

Dabei betrachte ich die Qualität der Schätzung gemeinsam mit dem Speicherbedarf und der gemessenen Ausführungszeit.

\end{sprechtext}

\folie{5}{NMC-Daten und Alterungsreferenz}{3:15--3:50}{0:35}{Einen Verlauf in der Grafik zeigen. Keine einzelnen Szenarien aufzählen.}

\begin{sprechtext}

Die erste Untersuchung basiert auf vierzehn NMC-Zellen in sieben Alterungsszenarien. Die Kurven zeigen, dass sich die Kapazität je nach Betriebsbedingungen unterschiedlich entwickelt. Einfluss haben unter anderem Temperatur, Ladezustand, Entladetiefe und Entladerate.

Der SOH wird hier über wiederholte Kapazitätsmessungen referenziert. Er beschreibt die verbleibende Kapazität im Verhältnis zur Nennkapazität. Eine gealterte Zelle kann also vollständig geladen sein und trotzdem weniger Ladung speichern als eine neue.

Diesen Datensatz verwende ich für die Entwicklung der SOH-Modelle.

\end{sprechtext}

\folie{6}{LFP-Daten für Benchmark und Optimierung}{3:50--4:25}{0:35}{Den Wechsel des Datensatzes deutlich machen.}

\begin{sprechtext}

Für den Robustheitsvergleich und die spätere Embedded-Optimierung verwende ich einen zweiten Datensatz mit fünfzehn LFP-Zellen. Hier wurden Lade- und Entladerate sowie Entladetiefe systematisch variiert, bei einer Umgebungstemperatur von fünfundzwanzig Grad.

Auch hier bilden Kapazitätsmessungen die Referenz für den Alterungszustand. Die verschiedenen Verläufe ermöglichen es, Modelle über unterschiedliche Belastungen und Alterungsstände hinweg zu untersuchen.

Wichtig für die Einordnung der folgenden Ergebnisse ist also: Der erste Modellteil verwendet NMC-Daten, die beiden weiteren Untersuchungen verwenden LFP-Daten.

\end{sprechtext}

\folie{7}{BMS-Plattform und Benchmark-Hardware}{4:25--5:05}{0:40}{Die beiden markierten Bausteine zeigen. H755 und H753 unterscheiden.}

\begin{sprechtext}

Die entwickelte BMS-Plattform verbindet Messung, Überwachung und eigene Zustandsmodelle. Der hier markierte BMS-Baustein misst Zellspannungen und Temperaturen und ermöglicht den passiven Ladungsausgleich.

Auf dem STM32H755 sind die Aufgaben auf zwei Kerne verteilt. Der Cortex-M4 übernimmt die hardwarenahen Aufgaben wie Messwerterfassung und Kommunikation. Der Cortex-M7 ist für die Zustandsmodelle vorgesehen.

Die später gezeigten Speicher- und Laufzeitmessungen wurden auf einem STM32H753 durchgeführt. Die eigene BMS-Platine zeigt damit die Integrationsumgebung, während der H753 die konkrete Plattform für den Hardwarebenchmark ist.

\end{sprechtext}

\folie{8}{Erster Untersuchungsteil: Eingabedarstellung}{5:05--5:35}{0:30}{Den Schwerpunkt auf Feature Engineering legen, nicht auf Gate-Gleichungen.}

\begin{sprechtext}

Der erste Untersuchungsteil setzt vor der Wahl einer aufwendigen Architektur an. Die Frage ist: Wie viel zeitabhängiges Batterieverhalten lässt sich bereits über eine geeignete Eingabedarstellung erfassen?

Als Ausgangspunkt dient ein MLP, also ein vollständig verbundenes neuronales Netz. Es besitzt keinen internen zeitlichen Zustand. Deshalb untersuche ich, wie abgeleitete Merkmale, Zeitauflösung und eine explizit zusammengestellte Historie die SOH-Schätzung beeinflussen.

\end{sprechtext}

\folie{9}{Wie eine Lag Sequence entsteht}{5:35--6:30}{0:55}{In der Grafik von den einzelnen Messgrößen zum gemeinsamen Eingangsvektor gehen.}

\begin{sprechtext}

Die Historie wird über sogenannte Lag Sequences bereitgestellt. Dafür fasse ich die Messwerte zunächst zu Zeitintervallen zusammen. Aus jedem Intervall entstehen Merkmale aus Spannung, Temperatur und kumuliertem Lade- beziehungsweise Entladestrom.

Mehrere aufeinanderfolgende Intervalle werden anschließend zu einem einzigen Eingangsvektor zusammengefügt. Das MLP bekommt damit sowohl aktuelle als auch vergangene Informationen. Der zugehörige Zielwert ist der SOH am Ende des Fensters.

Die zeitliche Abdeckung ergibt sich aus der Anzahl der Schritte und ihrer Dauer. Sechzehn Schritte mit jeweils zehn Minuten ergeben hier eine Historie von hundertsechzig Minuten. Die zeitliche Information steckt damit in den Eingängen; sie wird nicht durch einen rekurrenten Zustand im Netz gespeichert.

\end{sprechtext}

\folie{10}{Lernen, Modellwahl und unabhängiger Test}{6:30--7:20}{0:50}{Zelltrennung zuerst erklären, danach MSE und MAE unterscheiden.}

\begin{sprechtext}

Für Training und Validierung werden zehn NMC-Zellen verwendet. Vier weitere Zellen bleiben für den abschließenden Test reserviert. Zwanzig Prozent der Sequenzen aus den zehn Entwicklungszellen dienen der Validierung und der Auswahl der Hyperparameter.

Die Eingänge werden einheitlich auf den Bereich zwischen null und eins skaliert. Beim Training werden die Gewichte so angepasst, dass der mittlere quadratische Fehler, also der MSE, sinkt. Die Modellsuche variiert unter anderem Tiefe, Breite und zeitliche Eingabekonfiguration.

Die abschließende Evaluation erfolgt auf den vier separaten Testzellen mit MAE und MSE. Dadurch wird sichtbar, wie gut sich das Modell auf Zellen überträgt, die nicht zum eigentlichen Training gehören.

\end{sprechtext}

\folie{11}{Zeitauflösung und Historienlänge}{7:20--8:00}{0:40}{Eine passende Region der Fehlermatrix zeigen, nicht alle Kombinationen besprechen.}

\begin{sprechtext}

Diese Auswertung zeigt den Einfluss von Zeitauflösung und Sequenzlänge. Sehr grobe Intervalle glätten zwar Messwerte, können aber relevante Lade- und Entladedynamik verlieren.

Bei einer feinen Auflösung deckt dieselbe Anzahl von Schritten dagegen weniger Zeit ab. Wenn ich gleichzeitig eine lange Historie behalten möchte, wächst der Eingangsvektor.

In der untersuchten Konfiguration ergibt sich mit sechzehn Schritten und zehn Minuten pro Schritt ein geeigneter Kompromiss. Entscheidend ist also das Zusammenspiel von zeitlicher Abdeckung und Detailgrad. Mehr Historie oder feinere Daten führen nicht automatisch zu einer besseren Schätzung.

\end{sprechtext}

\folie{12}{SOH-Ergebnisse und Grenze der Übertragbarkeit}{8:00--8:50}{0:50}{Zuerst gute Verläufe, dann Zelle 9 als Gegenbeispiel zeigen.}

\begin{sprechtext}

Bei den hier gezeigten Zellen eins und sieben folgt die SOH-Schätzung der Kapazitätsreferenz recht gut. Die mittleren absoluten Fehler liegen je nach gezeigtem Verlauf ungefähr zwischen einem Viertel und einem halben Prozentpunkt.

Zelle neun zeigt die Grenze deutlicher. Hier treten systematische und lokale Abweichungen auf, mit einem MAE von rund zweieinhalb Prozentpunkten.

Die Untersuchung zeigt damit, dass eine Feedforward-Architektur mit geeigneter Historie zeitabhängige Alterung abbilden kann. Sie zeigt aber auch, dass die Übertragung auf andere Alterungsverläufe nicht selbstverständlich ist. Diese Frage nach der Zuverlässigkeit führt zum nächsten Teil: Wie verhalten sich Schätzer unter gezielt veränderten Bedingungen?

\end{sprechtext}

\folie{13}{Zweiter Untersuchungsteil: Robustheit}{8:50--9:20}{0:30}{Den Wechsel von SOH-Modellentwicklung zu SOC-Schätzervergleich klar aussprechen.}

\begin{sprechtext}

Im zweiten Teil geht es um einen Vergleich von SOC-Schätzern. Ich untersuche vier konkrete Vertreter unterschiedlicher Modellklassen unter denselben Störungen.

Neben dem normalen Schätzfehler betrachte ich die Reaktion auf Messfehler und die Erholung nach einer falschen Initialisierung. Ergänzend werden die Implementierungen auf dem Mikrocontroller bewertet.

Ich stelle zunächst kurz vor, welche Informationen die vier Verfahren jeweils nutzen. Genau daraus lassen sich später viele ihrer Stärken und Schwächen erklären.

\end{sprechtext}

\folie{14}{DM: Stromintegration mit fester Kapazität}{9:20--9:50}{0:30}{Den Integrationspfad einmal von links nach rechts zeigen.}

\begin{sprechtext}

Der erste Ansatz ist Direct Measurement, hier als Stromintegration mit fester Nennkapazität umgesetzt. Aus Strom und verstrichener Zeit wird die umgesetzte Ladung berechnet und auf die Kapazität bezogen.

Ein erkannter Volladezustand dient als Anker. Zwischen solchen Ankern gibt es keine kontinuierliche Korrektur über die Zellspannung.

Das Verfahren ist transparent und benötigt sehr wenig Rechenleistung. Es reagiert aber empfindlich auf einen falschen Anfangszustand, Strommessfehler und eine durch Alterung veränderte Kapazität.

\end{sprechtext}

\folie{15}{HDM: Kapazität mit SOH anpassen}{9:50--10:30}{0:40}{Nur den neu hinzugekommenen SOH-Pfad erklären.}

\begin{sprechtext}

HDM erweitert denselben Integrationskern um eine Anpassung an die Alterung. Ein LSTM schätzt aus stündlich aufbereiteten Merkmalen den SOH. Daraus wird die effektive Kapazität als Nennkapazität mal SOH bestimmt.

Das LSTM kann zeitliche Information in seinen Zuständen weitertragen. Die Gates steuern dabei, was erhalten, ergänzt und ausgegeben wird.

Der zusätzliche Nutzen liegt hier in der Kapazitätsanpassung. Der SOC entsteht weiterhin aus der Ladungsbilanz. Deshalb beseitigt dieser Ansatz weder einen systematischen Stromfehler noch automatisch eine falsche SOC-Initialisierung. Außerdem geht die Unsicherheit der SOH-Schätzung in die Kapazität ein.

\end{sprechtext}

\folie{16}{HECM: Physikmodell mit Spannungsrückkopplung}{10:30--11:15}{0:45}{Vorhersage und Korrektur am Blockbild erklären; Formeln bleiben im Anhang.}

\begin{sprechtext}

HECM ergänzt eine physikalische Beschreibung des Batterieverhaltens. Das Zwei-RC-Ersatzschaltbild bildet den ohmschen Anteil und zwei verzögerte Spannungsanteile ab. Der SOH passt die Kapazität und die entsprechenden Kennfelder an.

Im ersten Schritt schreibt das Modell den Zustand mit dem Strom fort und berechnet eine erwartete Klemmenspannung. Anschließend vergleicht der Extended Kalman Filter diese Vorhersage mit der gemessenen Spannung.

Aus der Abweichung entsteht eine Korrektur für SOC und RC-Zustände. Wie stark korrigiert wird, hängt von den angenommenen Unsicherheiten ab. Diese Spannungsrückkopplung ist der wesentliche zusätzliche Informationsweg gegenüber den beiden Stromzählern.

\end{sprechtext}

\folie{17}{DD: GRU mit MLP-Head}{11:15--12:00}{0:45}{Die Signalhistorie, die GRU und den skalaren Ausgang verbinden.}

\begin{sprechtext}

Der datengetriebene Ansatz verwendet eine GRU mit siebenundsechzig rekurrenten Einheiten und einem nachgeschalteten MLP. Als Eingänge erhält er Spannung, Strom, Temperatur, kumulierte Ladung, dynamische Merkmale, Zeitschritt und SOH.

Die GRU verarbeitet die zeitlichen Zusammenhänge in einem Hidden State. Das MLP bildet diesen Merkmalsvektor anschließend auf eine einzelne SOC-Schätzung ab.

Im Robustheitsbenchmark wird für jede Ausgabe ein Fenster mit zweitausendvierundzwanzig Messschritten neu verarbeitet. Dieser Ausführungsmodus ist für die Interpretation wichtig: Ein später untersuchter kontinuierlicher Betrieb mit weitergetragenem Zustand ist nicht automatisch dasselbe Verhalten.

\end{sprechtext}

\folie{18}{Störung der Initialisierung}{12:00--12:25}{0:25}{Auf den Startfehler im Schema zeigen.}

\begin{sprechtext}

Die erste Störungsgruppe betrifft den Anfangszustand. Nach einem Neustart oder einer längeren Lagerung kann dem BMS eine verlässliche SOC-Historie fehlen.

Deshalb starte ich die Schätzer gezielt mit einem falschen SOC. Untersucht wird anschließend, ob und wie schnell sie sich erholen. Bei LFP ist das besonders interessant, weil die flache Leerlaufspannungskennlinie in weiten Bereichen nur wenig Information über den SOC liefert.

\end{sprechtext}

\folie{19}{Störungen der Messwerte}{12:25--12:55}{0:30}{Gain, Offset und Rauschen als unterschiedliche Fehlerarten nennen.}

\begin{sprechtext}

Die zweite Gruppe umfasst Störungen direkt in den Eingangssignalen. Dazu gehören beispielsweise ein Verstärkungsfehler, ein Offset oder zusätzliches Rauschen.

Ein Stromfehler verändert die Ladungsbilanz. Ein Spannungsfehler kann das Messfeedback des Kalman Filters oder die Eingänge des neuronalen Netzes verfälschen.

Mir geht es dabei darum, die unterschiedlichen Reaktionsmechanismen unter vergleichbaren Bedingungen sichtbar zu machen. Alle untersuchten Schätzer erhalten deshalb dieselbe jeweilige Störung.

\end{sprechtext}

\folie{20}{Störungen der Signalübertragung}{12:55--13:20}{0:25}{Vom Sensor auf die Übertragung zum Algorithmus wechseln.}

\begin{sprechtext}

Die dritte Gruppe betrifft die Verfügbarkeit und zeitliche Zuordnung der Signale. Beispiele sind Paketverluste, Kommunikationsausfälle und verzögerte Messwerte.

Die reale Batterie arbeitet währenddessen weiter. Ein fehlender Datenpunkt bedeutet also nicht, dass auch kein Strom geflossen ist.

Solche Fehler können sowohl die aufsummierte Ladung als auch die zeitliche Historie eines neuronalen Modells verändern. Deshalb werden sie zusätzlich zu den reinen Sensorfehlern untersucht.

\end{sprechtext}

\folie{21}{Wie die Störungen ausgewertet werden}{13:20--14:00}{0:40}{Den Pfad vom gestörten Signal bis zu den neu berechneten Features verfolgen.}

\begin{sprechtext}

Für einen fairen Vergleich müssen sich die Störungen auch durch die Datenaufbereitung fortsetzen. Wenn beispielsweise der Strom gestört ist, berechne ich daraus auch die kumulierte Ladung und weitere abhängige Merkmale neu. Sonst könnte das Modell über unveränderte Features noch auf saubere Informationen zugreifen.

Anschließend vergleiche ich jeden gestörten Lauf mit seinem passenden ungestörten Lauf, also bei derselben Zelle, demselben Fenster und derselben Zufallsrealisierung.

Die Fehlerzunahme beschreibt die zusätzliche Belastung durch die Störung. Bei der Recovery interessiert mich dagegen, wann der falsch initialisierte Lauf wieder zum korrekt initialisierten Lauf zurückfindet.

\end{sprechtext}

\folie{22}{Nominale Genauigkeit als Ausgangspunkt}{14:00--14:55}{0:55}{Zuerst den linken MAE-Plot erklären. Punkte und Balken kurz zuordnen.}

\begin{sprechtext}

Zunächst sehen wir die Ergebnisse ohne zusätzliche Störung. Die Testbasis besteht aus sechs unabhängigen LFP-Zellen und sechzehn repräsentativen Tagesfenstern. Jede Zelle erhält bei der Zusammenfassung dasselbe Gewicht.

Die Punkte zeigen die einzelnen Fensterwerte, die Balken das Makromittel und die Fehlerbalken das Konfidenzintervall. DD erreicht hier den niedrigsten mittleren MAE mit ungefähr zweieinhalb Prozentpunkten. Zugleich liegen seine Fensterwerte enger zusammen.

Bei den anderen Verfahren ist die Streuung größer. Beim einfachen Stromzähler bleibt die Kapazität trotz Alterung fest. Die hybriden Verfahren berücksichtigen Alterung, erhalten dadurch aber auch einen zusätzlichen Einfluss der SOH-Schätzung. Diese nominale Rangfolge ist der Ausgangspunkt für die anschließenden Störtests.

\end{sprechtext}

\folie{23}{Reaktion auf Strom-Gain-Fehler}{14:55--15:40}{0:45}{Die zunehmende Fehlerkurve und dann den Vergleich der vier Modelle zeigen.}

\begin{sprechtext}

Hier wird der gemessene Strom mit einem systematischen Faktor verändert. Man kann sich das als Fehlkalibrierung des Sensors vorstellen.

DM und HDM zeigen den stärksten Fehleranstieg. Das ist plausibel, weil der falsch skalierte Strom fortlaufend integriert wird. Die SOH-Anpassung von HDM korrigiert diesen Messfehler nicht.

HECM kann einen Teil der Abweichung über die Spannung zurückführen. DD zeigt in diesem Vergleich den geringsten Anstieg. Eine plausible Erklärung ist, dass es mehrere Messgrößen und deren zeitlichen Zusammenhang gemeinsam verarbeitet. Aus diesem Ergebnis allein lässt sich aber noch nicht beweisen, welche einzelnen Merkmale dafür verantwortlich sind.

\end{sprechtext}

\folie{24}{Recovery: Rückkehr und dauerhafte Erholung}{15:40--16:40}{1:00}{Erst den Zeitverlauf, dann die beiden Recovery-Balken unterscheiden.}

\begin{sprechtext}

Bei einer falschen Initialisierung wird eine andere Stärke sichtbar. Der obere Verlauf zeigt einen anfänglichen SOC-Fehler von zehn Prozentpunkten. Die untere Darstellung betrachtet die Abweichung vom jeweils korrekt initialisierten Lauf.

Dabei unterscheide ich den ersten Eintritt in das Erholungsband von einer dauerhaften Rückkehr. Ein Modell kann das Band zunächst erreichen und es später wieder verlassen.

HECM zeigt hier die schnellste persistente Erholung. Die kontinuierliche Spannungsrückkopplung ermöglicht eine erneute Verankerung des SOC. DD kann sich ebenfalls erholen, zeigt aber im Beispiel zwischenzeitliche Rückfälle. DM und HDM bleiben ohne passenden Anker überwiegend verschoben.

Damit sehen wir: Der Schätzer mit dem kleinsten nominalen Fehler muss nicht in jeder anderen Bewertung ebenfalls vorne liegen.

\end{sprechtext}

\folie{25}{Spannungsspitzen und lokale Fehler}{16:40--17:30}{0:50}{Den kurzen DD-Ausschlag zeigen und anschließend auf den langen Auswertungszeitraum verweisen.}

\begin{sprechtext}

Ein kurzer Spannungsimpuls zeigt wiederum eine andere Schwachstelle. DD reagiert am Ereignis mit einem deutlichen SOC-Ausschlag, der danach rasch abklingt.

Das passt zur direkten Verarbeitung von Spannung und Spannungsänderung als Eingangsmerkmale. Die genaue Aufteilung dieser Einflüsse müsste man allerdings mit zusätzlichen Untersuchungen prüfen.

Über ein ganzes Tagesfenster verändert der kurze Ausschlag den mittleren Fehler kaum. Für eine Anwendung kann eine solche lokale Reaktion trotzdem relevant sein.

Deshalb reicht es nicht, nur den MAE über die gesamte Messreihe zu betrachten. Die zeitliche Reaktion auf einzelne Ereignisse ergänzt die globale Fehlerbewertung.

\end{sprechtext}

\folie{26}{Was der Robustheitsvergleich insgesamt zeigt}{17:30--18:10}{0:40}{DD-Gesamtscore und HECM-Recovery als unterschiedliche Stärken benennen.}

\begin{sprechtext}

In der zusammengeführten Bewertung erzielt DD bei den drei betrachteten Prioritätsprofilen den höchsten Gesamtscore. Die Einzeluntersuchungen zeigen gleichzeitig, wo dieses Ergebnis differenziert werden muss.

HECM ist besonders stark bei der dauerhaften Erholung nach einem Initialisierungsfehler. DD ist bei vielen Störungen robust, kann aber auf lokale Spannungsspitzen empfindlich reagieren.

Die Scores sind relative Vergleichsgrößen innerhalb dieser Untersuchung. Für die Auswahl im BMS sind deshalb sowohl das Gesamtbild als auch die konkreten Fehlermechanismen wichtig. Nach der Zuverlässigkeit folgt jetzt die Frage, welchen Aufwand die neuronalen Modelle auf der Hardware verursachen.

\end{sprechtext}

\folie{27}{Dritter Untersuchungsteil: Embedded-Optimierung}{18:10--18:45}{0:35}{Pruning und Quantisierung als zwei getrennte Wege einführen.}

\begin{sprechtext}

Im dritten Teil untersuche ich eigenständige LSTM-Modelle für SOC und SOH. Diese Modelle sind die Ausgangspunkte für zwei getrennte Optimierungswege: strukturiertes Pruning und eine INT8-Quantisierung der Gewichte.

Beim Pruning wird die Netzstruktur verkleinert. Bei der Quantisierung werden ausgewählte Gewichte mit weniger Bits gespeichert.

Beide Varianten vergleiche ich mit dem jeweiligen Basismodell auf dem PC und auf dem STM32. Entscheidend ist, welche Einsparungen tatsächlich entstehen und wie sich gleichzeitig die Schätzqualität verändert.

\end{sprechtext}

\folie{28}{Parameter im LSTM-Kern}{18:45--19:35}{0:50}{Input-Gewichte, rekurrente Gewichte und den quadratischen Anteil H² zeigen.}

\begin{sprechtext}

Zunächst schauen wir darauf, wo die Parameter liegen. Beide Modelle erhalten sechs Eingangsmerkmale. Das SOC-LSTM besitzt vierundsechzig Hidden Units, das SOH-LSTM hundertachtundzwanzig.

Für jede der vier Transformationen gibt es Gewichte für den aktuellen Eingang und für den vorherigen Hidden State. Die Eingangsmatrix enthält H mal D Werte. Die rekurrente Matrix enthält H mal H Werte.

In der verwendeten Implementierung kommen je Transformation zwei Bias-Vektoren hinzu. Daraus ergibt sich die hier angegebene Parameterzahl.

Wichtig ist besonders der quadratische Anteil: Wenn die rekurrente Breite wächst, nimmt die Zahl der Verbindungen innerhalb des LSTM sehr schnell zu.

\end{sprechtext}

\folie{29}{Parameter im MLP-Head}{19:35--20:05}{0:30}{Am rechten Teil der Architektur bleiben.}

\begin{sprechtext}

Der MLP-Head macht aus dem Hidden State die eigentliche SOC- oder SOH-Schätzung. Die erste Schicht verbindet die H LSTM-Ausgänge mit M Neuronen. Danach folgt ein einzelner Ausgang.

Damit entstehen H mal M Gewichte, zwei mal M weitere Parameter für Bias und Ausgangsgewichte sowie ein letzter Bias.

Für die beiden Basismodelle entspricht die MLP-Breite jeweils der LSTM-Breite. Der Head wird bei der folgenden Gesamtzahl mit berücksichtigt.

\end{sprechtext}

\folie{30}{Gesamtzahl und Ansatzpunkt der Kompression}{20:05--20:35}{0:30}{Rund 81 Prozent im LSTM hervorheben, nicht jede Tabellenzahl vorlesen.}

\begin{sprechtext}

Zusammen ergibt das rund dreiundzwanzigtausend Parameter für das SOC-Modell und sechsundachtzigtausend für das SOH-Modell.

Bei beiden liegt ungefähr einundachtzig Prozent der Parameter im LSTM-Kern. Dort befindet sich außerdem der quadratische Anteil der rekurrenten Verbindungen.

Deshalb setzt die strukturelle Verkleinerung an den Hidden Units des LSTM an. Wenn ich dort eine Einheit entferne, betrifft das mehrere gekoppelte Matrizen und zugleich die Anbindung an den MLP-Head.

\end{sprechtext}

\folie{31}{Vom trainierten Modell bis zum Hardwarevergleich}{20:35--21:15}{0:40}{Die beiden Variantenpfade und den gemeinsamen Vergleichspunkt zeigen.}

\begin{sprechtext}

Der Versuchsaufbau beginnt bei denselben trainierten SOC- und SOH-Basismodellen. Daraus entstehen entweder eine geprunte Variante mit anschließendem Fine-Tuning oder eine quantisierte Variante.

Die Parameter werden in die C-Implementierung übertragen. Anschließend erhalten PC und Mikrocontroller denselben Datenstrom, damit sich ihre Ausgaben vergleichen lassen.

Zusätzlich messe ich Flash, RAM und Ausführungszeit. Damit wird nicht nur die Größe einer gespeicherten Modelldatei bewertet, sondern der Aufwand in der tatsächlichen Firmware.

Pruning und Quantisierung werden hier getrennt verglichen. Eine kombinierte Variante ist nicht Teil dieses Experiments.

\end{sprechtext}

\folie{32}{Warum strukturiertes Pruning?}{21:15--21:55}{0:40}{Die einzelnen Nullen mit den entfernten Zeilen und Spalten vergleichen.}

\begin{sprechtext}

Der Unterschied zwischen unstrukturiertem und strukturiertem Pruning ist für diese Implementierung entscheidend.

Wenn ich einzelne Gewichte auf null setze, bleiben die Dimensionen einer dicht gespeicherten Matrix zunächst gleich. Eine normale dichte Rechenschleife läuft weiterhin über diese Positionen.

Beim strukturierten Pruning entferne ich ganze Kanäle und die dazugehörigen Zeilen und Spalten. Dadurch entstehen tatsächlich kleinere Matrizen.

Diese kleineren Tensoren kann der vorhandene FP32-Kernel direkt verarbeiten. Für den gewählten Mikrocontrolleraufbau ist das der praktische Vorteil: Die verkleinerte Struktur lässt sich ohne spezielle Sparse-Verarbeitung nutzen.

\end{sprechtext}

\folie{33}{L2-Ranking und konsistente Kanalentfernung}{21:55--23:00}{1:05}{Score, Sortierung und gekoppelte Matrizen nacheinander zeigen.}

\begin{sprechtext}

Für die Auswahl bekommt jede Hidden Unit einen Score aus den L2-Normen ihrer Gewichtszeilen. Dabei werden die Eingangs- und rekurrenten Gewichte über alle vier LSTM-Blöcke berücksichtigt.

Die L2-Norm beschreibt hier die Größe der Gewichte. Sie ist kein Vorhersagefehler. Ich berechne also nicht für jede Zeile einen MSE, sondern sortiere die Kanäle nach diesem Gewichtsscore und entferne den vorgegebenen Anteil mit den kleinsten Scores.

Danach muss dieselbe Auswahl an allen abhängigen Stellen umgesetzt werden: in den Gates, den rekurrenten Zeilen und Spalten, den Bias-Vektoren, den Zuständen und den Eingangsspalten des MLP.

So sinkt die Hidden Size beim SOC von vierundsechzig auf fünfundvierzig und beim SOH von hundertachtundzwanzig auf neunzig. Anschließend werden die verbleibenden Gewichte nachtrainiert, damit sich das kleinere Netz anpassen kann.

\end{sprechtext}

\folie{34}{Was Weight-only INT8 hier bedeutet}{23:00--24:15}{1:15}{Eine Zeile der Gewichtsmatrix, ihre Skala und die Rekonstruktion erläutern.}

\begin{sprechtext}

Bei der Quantisierung bleibt die Netzstruktur erhalten. Stattdessen werden die LSTM-Gewichte zeilenweise symmetrisch auf INT8 abgebildet.

Für jede Zeile bestimme ich zunächst den größten Gewichts-Betrag. Daraus entsteht die Skala, indem ich durch hundertsiebenundzwanzig teile. Dann teile ich jedes Gewicht durch diese Skala, runde und begrenze den Wert auf den erlaubten Integerbereich. Gespeichert werden die ganzzahligen Werte und eine Skala pro Zeile.

Ein solcher Gewichtswert benötigt damit ein Byte statt vier. Die Skalen benötigen zusätzlichen Speicher, und andere Teile des Modells bleiben unverändert.

Zur Laufzeit wird aus Integerwert mal Skala wieder ein FP32-Gewicht rekonstruiert. Auch Bias, Zustände und MLP bleiben in FP32. Das ist also eine kompaktere Gewichtsspeicherung bei weiterhin gleitkommabasierter Berechnung. Die Anzahl der Rechenverbindungen bleibt gleich. Deshalb muss der Flash-Gewinn nicht mit einem Laufzeitgewinn einhergehen.

\end{sprechtext}

\folie{35}{Schätzfehler nach der Kompression}{24:15--25:05}{0:50}{SOC links und SOH rechts getrennt bewerten.}

\begin{sprechtext}

Bei SOC bleiben die Fehler in einer ähnlichen Größenordnung. Der MAE liegt bei ungefähr zwei Komma sieben Prozentpunkten für Base, zwei Komma drei für Pruned und zwei Komma acht für Quantized. Die geprunte Variante zeigt hier sogar eine geringere Streuung.

Bei SOH ist der Unterschied deutlicher. Base erreicht null Komma fünfundachtzig Prozentpunkte, die beiden komprimierten Varianten ungefähr eins Komma vier bis eins Komma fünf.

Diese SOH-Werte beziehen sich auf die gezeigte kausale Filterung. Insgesamt bleibt die SOC-Genauigkeit weitgehend erhalten, während die Ressourceneinsparung beim SOH mit höheren Fehlern verbunden ist. Kompression muss daher für jede Aufgabe separat bewertet werden.

\end{sprechtext}

\folie{36}{Flash- und RAM-Einsparung}{25:05--25:55}{0:50}{Zuerst den Flash-Plot, anschließend die kleinere RAM-Ersparnis erklären.}

\begin{sprechtext}

Der Flash-Bedarf sinkt bei beiden Optimierungswegen deutlich. Beim SOC-Modell geht er von etwa hundertsechs KiB auf siebenundsechzig KiB durch Pruning und auf vierundfünfzig KiB durch Quantisierung zurück.

Relativ spart die Quantisierung hier rund fünfzig Prozent Flash beim SOC und sechsundfünfzig Prozent beim SOH. Pruning erreicht ungefähr siebenunddreißig beziehungsweise dreiundvierzig Prozent.

Die RAM-Einsparung fällt mit ungefähr achtzehn bis dreiundzwanzig Prozent kleiner aus. Im RAM liegen neben Parametern auch Zustände, Aktivierungen, Puffer und Stack. Bei der quantisierten Variante bleiben viele dieser Größen FP32. Deshalb folgt der gesamte Hardwarebedarf nicht einfach der Bitbreite der gespeicherten Gewichte.

\end{sprechtext}

\folie{37}{Pruning: Kernelzeit und Host-Latenz}{25:55--26:50}{0:55}{Die beiden Zeitmessungen ausdrücklich auseinanderhalten.}

\begin{sprechtext}

Bei der Laufzeit unterscheide ich zwei Messgrößen. Die Kernelzeit misst die Berechnung auf dem STM32. Die Host-Latenz umfasst zusätzlich die Übertragung und den Ablauf zwischen PC und Controller.

Pruning verkürzt die Kernelzeit um ungefähr dreiundvierzig Prozent beim SOC und vierundvierzig Prozent beim SOH. Die kleineren Matrizen führen also auch in der Umsetzung zu einem deutlichen Rechenzeitgewinn.

Bei der Host-Latenz fällt der Gewinn unterschiedlich aus: ungefähr vier Prozent beim SOC und neunundzwanzig Prozent beim SOH. Beim kurzen SOC-Kernel beanspruchen Übertragung und Scheduling einen großen Anteil. Eine schnellere Netzberechnung verkürzt daher nicht automatisch den gesamten Ablauf im selben Verhältnis.

\end{sprechtext}

\folie{38}{Quantisierung: Flash-Gewinn und Zeitaufwand}{26:50--27:40}{0:50}{Das Ergebnis auf die konkrete gemischte INT8/FP32-Implementierung beziehen.}

\begin{sprechtext}

Die quantisierte Variante zeigt das Gegenstück. Sie benötigt weniger Flash, aber in diesem Referenzkernel mehr Zeit. Beim SOC ist die Kernelzeit ungefähr fünfmal so lang wie bei Base. Beim SOH steigt sie um etwa achtundzwanzig Prozent.

Die Gewichte werden während der FP32-Berechnung rekonstruiert. Das verursacht zusätzlichen Aufwand; auch Speicherzugriffe und Schleifen können zum gemessenen Ergebnis beitragen.

Damit ist die Quantisierung in dieser Umsetzung ein Speicherhebel. Für eine Beschleunigung müsste man den Kernel gezielt untersuchen und gegebenenfalls andere Rechenpfade nutzen. Das Ergebnis beschreibt diese Implementierung und lässt sich nicht pauschal auf jede INT8-Ausführung übertragen.

\end{sprechtext}

\folie{39}{Antworten auf die Forschungsfragen}{27:40--28:50}{1:10}{Die drei Ergebnisse noch einmal in derselben Reihenfolge wie auf Folie 3 verbinden.}

\begin{sprechtext}

Damit komme ich zurück zu den drei Forschungsfragen.

Erstens zeigt die MLP-Untersuchung, dass sich zeitabhängige SOH-Verläufe bereits mit einer Feedforward-Architektur beschreiben lassen, wenn Belastungshistorie und zeitlicher Kontext gezielt in den Eingängen enthalten sind. Die Datenrepräsentation ist damit ein wichtiger Teil des Modellentwurfs.

Zweitens zeigt der Benchmark, dass nominale Genauigkeit allein für die Bewertung eines BMS-Schätzers nicht ausreicht. DD erzielt insgesamt starke Ergebnisse. Die Recovery von HECM und die lokale Empfindlichkeit von DD machen aber deutlich, warum Störtests und die Verwaltung der Modellzustände dazugehören.

Drittens ist neuronale SOC- und SOH-Schätzung auf dem Mikrocontroller realisierbar. Strukturiertes Pruning spart Speicher und Rechenzeit. Weight-only-Quantisierung spart vor allem Flash, mit einem zusätzlichen Zeitaufwand im verwendeten Kernel.

Der gemeinsame Beitrag ist damit die Bewertung vom Eingangssignal bis zur Firmware: Welche Information braucht das Modell, wie zuverlässig nutzt es diese Information und welche Ressourcen benötigt die Umsetzung?

\end{sprechtext}

\folie{40}{Abschluss}{28:50--29:00}{0:10}{Nach dem letzten Satz eine kurze Pause lassen und zur Kommission schauen.}

\begin{sprechtext}

Vielen Dank für Ihre Aufmerksamkeit. Ich freue mich auf Ihre Fragen und die anschließende Diskussion.

\end{sprechtext}

%TC:ignore
\clearpage
\pdfbookmark[0]{Vorbereitung für Rückfragen}{rueckfragen}
{\Large\bfseries Vorbereitung für Rückfragen\par}
{\small\color{quiet}Ab hier kein Bestandteil des 30-Minuten-Vortrags.\par}

\textbf{Schneller Zugriff auf den Anhang}\par
\begin{tabularx}{\linewidth}{@{}>{\raggedright\arraybackslash}Xl@{}}
\toprule
Thema & PowerPoint-Folien \\
\midrule
LSTM- und GRU-Zelle mit Gleichungen & 44--45 \\
Trainings-, Validierungs- und Testtrennung & 46 \\
Fensterbetrieb und kontinuierlicher Zustand & 49 \\
Konkrete SOC-/SOH-LSTM-Architektur aus Kapitel 7 & 50--51 \\
Parameter, MACs und gekoppelte Pruning-Indizes & 52--53 \\
MLP-Gleichung und konkrete Architektur aus Kapitel 5 & 58--59 \\
GRU und gemeinsames SOH-LSTM aus Kapitel 6 & 60--61 \\
Hardwarebedarf der SOC-Verfahren & 65 \\
Signifikanz, Konfidenzintervalle und Versuchsbedingungen & 66--68 \\
Referenzwerte, Testfenster und Hyperparameterwahl & 69--72 \\
HECM-Übersicht und ECM in Übungsnotation & 73--76 \\
Kalmanfilter: Q, P, R, Prediction, Gain und Update & 77--79 \\
Fehlermaße, Paarvergleiche und Recovery & 80--83 \\
INT8-Rekonstruktion und kausale SOH-Filterung & 84--85 \\
Scores, Rechenersparnis und Quantisierung & 86--88 \\
Auswahl der Benchmarkfenster & 89 \\
\bottomrule
\end{tabularx}

\vspace{4mm}
\textbf{Begriffe, die beim Sprechen sauber getrennt bleiben}\par
SOC und SOH beschreiben Ladestand und Kapazitätszustand. Ein Fehler von
0{,}01 bei einem auf 0 bis 1 normierten Zustand entspricht einem Prozentpunkt.
Die Scores des Robustheitsvergleichs sind keine Genauigkeitsprozente.

Der L2-Pruning-Score ordnet Kanäle nach ihren Gewichten. Der MSE bewertet
Vorhersagen anhand von Referenzwerten. Pruning und Quantisierung sind im
gezeigten Versuch getrennte Varianten. Weight-only INT8 bedeutet in dieser
Implementierung weiterhin FP32-Rechnung mit rekonstruierten Gewichten.

Beim Kalmanfilter beschreibt Q das Prozessrauschen, R das Messrauschen und
P die Unsicherheit der Zustandsschätzung. K gewichtet die Messkorrektur.
Auf den Übungsfolien wird vom Zustand bei k zum Zustand bei k+1 gerechnet.
Die Zeitfolie 76 ordnet gespeicherten Strom und neue Spannungsmessung dem
geprüften Codeablauf zu. Das 1RC-Lehrmodell ist von eurem 2RC-HECM zu unterscheiden.

\clearpage
\pdfbookmark[0]{Begrüßung und Generalprobe}{begrueszung-probe}
{\Large\bfseries Begrüßung und Generalprobe\par}
{\small\color{quiet}Persönliche Vorbereitung, nicht mitsprechen.\par}

\textbf{Professor Opferkuch.} Der Bezug zur Bachelorarbeit stammt aus deinen
Angaben: Betreuung an der TH Nürnberg zum Thema Mess- und Analyseverfahren
von Lithium-Ionen-Akkus. Für den Einstieg genügt dieser persönliche Zusammenhang.
Prüfstand, Dymola und die einzelnen Arbeitsschritte würden den Dank unnötig verlängern.

\textbf{Professor Li.} Weihan Li leitet laut dem offiziellen
\href{https://rwthcontacts.rwth-aachen.de/organization/ORG-74JAH/PER-NJBL3RJ}{RWTH-Organisationsverzeichnis}
die Juniorprofessur für \emph{Artificial Intelligence and Digitalization for
Batteries}. Der fachliche Bezug zu neuronalen Batteriemodellen ist damit direkt
gegeben. Für die Begrüßung reicht der Dank für die Begutachtung und die Anreise.
Eine persönliche Zusammenarbeit wird nicht behauptet. Recherche: 1. Oktober 2026.

\textbf{Vier Zeitkontrollen}\par
\begin{tabularx}{\linewidth}{@{}>{\raggedright\arraybackslash}Xr@{}}
\toprule
Zeitpunkt & Zielzeit nach Beginn \\
\midrule
Nach Folie 12: SOH-MLP abgeschlossen & 8:50 \\
Nach Folie 26: Robustheitsbenchmark abgeschlossen & 18:10 \\
Nach Folie 38: Hardwareergebnisse abgeschlossen & 27:40 \\
Nach Folie 40: Übergang zur Diskussion & 29:00 \\
\bottomrule
\end{tabularx}

\vspace{4mm}
\textbf{Wenn du etwa eine Minute zurückliegst}\par
Auf Folie 4 genügen die drei eigenen Beiträge in einem Satz.
Auf den Folien 18 bis 20 je ein Beispiel nennen und direkt weitergehen.
Auf Folie 30 reichen die beiden gerundeten Gesamtzahlen und der LSTM-Anteil.
Bei Folie 34 den Rechenweg in einem Satz zusammenfassen und die FP32-Ausführung
klar benennen. Die Schlussantworten auf Folie 39 bleiben erhalten.

\textbf{Generalprobe}\par
Den Text einmal laut mit den Folien sprechen. Die Uhr startet mit der Begrüßung,
nicht erst beim fachlichen Inhalt. An Diagrammen den Blick kurz vom Manuskript
lösen und tatsächlich auf die gemeinte Kurve oder den Block zeigen. Das
benötigt Zeit und sollte in der Probe mit enthalten sein. Wenn das natürliche
Tempo langsamer als der Plan ist, zuerst Wiederholungen und Zahlenreihen kürzen.

\textbf{Grundlage dieser Fassung}\par
Aktuelle Präsentation v1.12 mit 40 Hauptfolien und erweitertem Anhang;
Zahlen, Modellzuordnung und Bildreihenfolge folgen dieser Präsentation.
Die Begrüßung verwendet die von dir benannten Mitglieder der Kommission.
Die vollständigen Folienabbildungen dienen hier als Orientierung für die Probe.

%TC:endignore
\end{document}
